% Generated draft: requires mathematical and source review.
\documentclass[11pt]{article}
\usepackage[margin=25mm]{geometry}
\usepackage{fontspec}
\setmainfont{Latin Modern Roman}
\usepackage{amsmath,amssymb,mathtools,unicode-math}
\setmathfont{Latin Modern Math}
\usepackage{enumitem,needspace,xurl,hyperref}
\setlength{\parindent}{0pt}
\setlength{\parskip}{5pt}
\setlength{\emergencystretch}{3em}
\newcommand{\N}{\mathbb{N}}
\newcommand{\Z}{\mathbb{Z}}
\newcommand{\Q}{\mathbb{Q}}
\newcommand{\R}{\mathbb{R}}
\newcommand{\C}{\mathbb{C}}
\newcommand{\F}{\mathbb{F}}
\newcommand{\A}{\mathbb{A}}
\newcommand{\PP}{\mathbb{P}}
\newcommand{\E}{\mathbb{E}}
\newcommand{\eps}{\varepsilon}
\newcommand{\dd}{\,\mathrm{d}}
\begin{document}
\section*{UnsolvedMath ID: 701731}
\textbf{Source problem code:} catalogue entry
\subsection{Definitions and mathematical statement}
For subsets \(A,B\) of a group \(G\), write \(AB=\{ab:a\in A,\ b\in B\}\). A factorization \(G=AB\) has unique representations if the multiplication map
\[
A\times B\longrightarrow G,\qquad (a,b)\longmapsto ab
\]
is bijective.

The supplied question asks whether every countable group \(G\) admits infinite subsets \(A,B\subseteq G\) such that
\[
\forall g\in G\quad \exists!\,(a,b)\in A\times B\quad g=ab.
\]
The meaning of ``countable'' requires clarification: if it includes finite groups, the requirement that both subsets be infinite is incompatible with those cases. The possible intended restriction to countably infinite groups is not explicit in the supplied record and requires confirmation.

Using the conventional meaning, a group is \emph{topologizable} if it admits a nondiscrete Hausdorff topology making multiplication and inversion continuous. The source additionally states that the requested factorization exists whenever \(G\) is countable and topologizable. This positive-case assertion is supplied source information, not independently verified here.
\subsection{Short English statement}
The question asks whether every countable group can be expressed uniquely as products from two infinite subsets. The intended exclusion of finite groups needs confirmation. The source states that countable topologizable groups have such a factorization.
\subsection{Sources}
\begin{itemize}
\item Local UnsolvedMath registry, permanent ID \(701731\), exact source problem code \texttt{catalogue entry}, title \emph{Kourovka Notebook Problem 20.73}. The supplied registry labels the problem open; this status is not independently verified.
\item I. V. Protasov, proposer, 2022. The supplied attribution identifies Problem 20.73 in \emph{The Kourovka Notebook (Unsolved Problems in Group Theory)}, 20th issue (2022), and the 21st issue edition (2026), identified as \texttt{arXiv:1401.0300v48}. The collection description names E. I. Khukhro and V. D. Mazurov as editors. No separate citation for the stated topologizable case is supplied.
\end{itemize}
\end{document}
